\section*{Chapitre \ref{ch:g12:destruction-of-mountains} --- La destruction des montagnes}

\begin{solution}{exo:g12:destruction-of-mountains:1}
(a) le \omterm{def:g12:destruction-of-mountains:denudation}{taux de dénudation}~; (b) la \omterm{def:g12:destruction-of-mountains:load}{charge de fond}~; (c) la \omterm{def:g12:destruction-of-mountains:load}{charge
dissoute}~; (d) la \omterm{def:g12:destruction-of-mountains:diagenesis}{diagenèse}~; (e) le \omterm{def:g12:destruction-of-mountains:rebound}{rebond isostatique}.
\end{solution}

\begin{solution}{exo:g12:destruction-of-mountains:2}
(1) Un conglomérat. (2) Un grès. (3) Une argilite feuilletée. (4) Une
brèche.
\end{solution}

\begin{solution}{exo:g12:destruction-of-mountains:3}
\qty{1}{m/Ma}, c'est \qty{1000}{mm} en \num{1000000} d'années, soit
\qty{0.001}{mm} par an. (1) \qty{0.054}{mm} par an. (2) \qty{0.218}{mm}
par an. (3) $\qty{10}{mm}$ par an, c'est \qty{10}{m} en 1000 ans, soit
\qty{10000}{m/Ma}.
\end{solution}

\begin{solution}{exo:g12:destruction-of-mountains:4}
La charge solide représente \qty{85}{\percent} du total~; le total vaut donc $60/0.85
\approx 70.6$ millions de tonnes, et la \omterm{def:g12:destruction-of-mountains:load}{charge dissoute} environ $70.6 - 60 \approx
10.6$ millions de tonnes par an.
\end{solution}

\begin{solution}{exo:g12:destruction-of-mountains:5}
L'altération du granite, le transport par une rivière, le dépôt dans un
bassin, l'enfouissement sous des couches plus jeunes, la cimentation.
\end{solution}

\begin{solution}{exo:g12:destruction-of-mountains:6}
À \qty{30}{cm/s}, le grain de sable est arraché (au-delà de
\qty{20}{cm/s})~; le grain de silt reste en suspension. À \qty{5}{cm/s},
le grain de sable se dépose (en dessous de \qty{10}{cm/s}), mais le grain
de silt reste toujours en suspension, puisque \qty{0.1}{cm/s} suffit à
l'y maintenir.
\end{solution}

\begin{solution}{exo:g12:destruction-of-mountains:7}
Charge totale $4/0.85 \approx 4.71$ millions de tonnes, soit $\num{4.71e9}/2830 \approx
\qty{1.66e6}{m^3}$ de roche par an. Sur $\qty{5e10}{m^2}$~: $\num{3.3e-5}$~m
par an, environ \qty{33}{m/Ma}, en dessous de la médiane de
\qty{54}{m/Ma}.
\end{solution}

\begin{solution}{exo:g12:destruction-of-mountains:8}
(1) $1 - 2700/3300 \approx 0.182$. (2) Environ \qty{182}{m}. (3) Un peu
plus que les \qty{163}{m} du chapitre~: plus la croûte est légère par
rapport au manteau, moins la colonne remonte.
\end{solution}

\begin{solution}{exo:g12:destruction-of-mountains:9}
(1) \qty{8}{km} de roche. (2) $8000 \times 0.163 \approx \qty{1300}{m}$.
(3) $8000/200 = 40$ millions d'années.
\end{solution}

\begin{solution}{exo:g12:destruction-of-mountains:10}
(1) $794/43 \approx 18$ millions d'années. (2) La surface ne s'abaisse que
de $0.163$ fois l'épaisseur érodée~; la roche à enlever vaut donc $794/0.163
\approx \qty{4900}{m}$, ce qui prend environ $4900/43 \approx 110$ millions
d'années. (3) Les continents ne font pas que s'user~: la formation des
montagnes et le magmatisme des arcs continuent de les soulever et de leur
ajouter de la croûte nouvelle.
\end{solution}

\begin{solution}{exo:g12:destruction-of-mountains:11}
(1) $\num{1e14}/2830 \approx \qty{3.5e10}{m^3}$, soit environ
\qty{35}{km^3}. (2) $100/15.5 \approx 6.5$ ans.
\end{solution}

\begin{solution}{exo:g12:destruction-of-mountains:12}
$E = kh$ et $h = h_0 e^{-fkt}$ donnent $E(t) = kh_0 e^{-fkt}$. L'épaisseur
enlevée de $t = 0$ à $t \to \infty$ vaut $\int_0^{\infty} kh_0
e^{-fkt}\,\mathrm{d}t = kh_0/(fk) = h_0/f$. Pour abaisser la surface de
$h_0$, l'érosion doit enlever $h_0/f \approx 6.1\,h_0$, comme le dit la
proposition sur le rebond~: chaque mètre érodé n'abaisse la surface que de
$f$ mètre.
\end{solution}

\begin{solution}{exo:g12:destruction-of-mountains:13}
(1) $10 \times 1000 = \qty{10000}{mm}$, soit \qty{10}{m}. (2) Retirer une
charge de glace d'épaisseur $t$ revient à retirer une masse $\rho_{\text{glace}}\,t$
par unité de surface~; la colonne remonte jusqu'à déplacer d'autant moins
de manteau~: $2000 \times 900/3380 \approx \qty{530}{m}$. (3) Le sol
remonte de moins en moins vite à mesure qu'il approche de son nouvel
équilibre~: ce qui entraîne la remontée, c'est la part qui n'est pas encore
rattrapée.
\end{solution}

\begin{solution}{exo:g12:destruction-of-mountains:14}
(1) Avant 1900~: environ \qty{56}{m/Ma} (comme dans l'exemple). En
1987--2006~: $145/0.85 \approx 171$ millions de tonnes, $\num{1.71e11}/2830
\approx \qty{6.0e7}{m^3}$, sur $\qty{2.98e12}{m^2}$~: environ
\qty{20}{m/Ma}. (2) La première~: la charge plus récente est retenue par
les barrages et d'autres ouvrages. (3) Non~: la diminution vaut
$400 - 145 = 255$ millions de tonnes par an, dont les barrages expliquent
100 à 150, environ la moitié.
\end{solution}

\begin{solution}{exo:g12:destruction-of-mountains:15}
(1) Au moins \qty{6}{km}. (2) $310 - 252 = 58$ millions d'années au plus,
donc au moins $6000/58 \approx \qty{100}{m/Ma}$. (3) C'est entre la
médiane (\qty{54}{m/Ma}) et la moyenne (\qty{218}{m/Ma}) des bassins~: un
taux ordinaire pour une chaîne de montagnes.
\end{solution}

\begin{solution}{pb:g12:destruction-of-mountains:1}
\textbf{1.} $E = 0.209 \times 2000 \approx \qty{418}{m/Ma}$.

\textbf{2.} \qty{0.418}{mm} par an.

\textbf{3.} Plus vite que les deux~: c'est une chaîne jeune et pentue,
aussi rapide que les bassins sismiquement actifs, et même davantage.

\textbf{4.} Sa pente moyenne.

\textbf{5.} Une couche de $\qty{0.418e-3}{m}$ sur $\qty{1e6}{m^2}$
représente \qty{418}{m^3}, soit $418 \times 2830 \approx \qty{1.18e6}{kg}$~:
environ 1200 tonnes.

\textbf{6.} $2000/418 \approx 4.8$ millions d'années.

\textbf{7.} $h(t) = 2000\,e^{-0.209\,t}$ ($t$ en millions d'années).

\textbf{8.} $1/0.209 \approx 4.8$ millions d'années.

\textbf{9.} $\ln 2/0.209 \approx 3.3$ millions d'années.

\textbf{10.} $e^{-0.209 \times 11} = e^{-2.30} \approx 0.100$.

\textbf{11.} $1 - 2830/3380 \approx 0.163$.

\textbf{12.} La surface s'abaisse de $100 \times 0.163 \approx \qty{16}{m}$~;
la \omterm{def:g12:continental-crust:root}{racine} remonte d'environ \qty{84}{m}.

\textbf{13.} $h' = -0.163 \times 0.209\,h \approx -0.0341\,h$, donc $h(t) =
2000\,e^{-0.0341\,t}$.

\textbf{14.} $(2000 - 200)/0.163 \approx \qty{11000}{m}$ de roche.

\textbf{15.} Des roches qui se trouvaient à environ \qty{11}{km} de
profondeur~: des roches métamorphiques comme les schistes et les gneiss, et
des granites.

\textbf{16.} $1/0.0341 \approx 29$ millions d'années.

\textbf{17.} $29 \times \ln 10 \approx 67$ millions d'années.

\textbf{18.} $2000\,e^{-0.0341 \times 200} \approx 2000 \times 0.0011 \approx
\qty{2}{m}$.

\textbf{19.} $2000\,e^{-0.0341 \times 250} \approx 2000 \times 0.0002 \approx
\qty{0.4}{m}$~: une plaine. Un sommet de \qty{2037}{m} signifie que le
modèle seul ne convient pas~: la région a dû être soulevée de nouveau
depuis, ou s'éroder plus lentement que le modèle ne le suppose.

\textbf{20.} L'altitude moyenne est divisée par deux en $29 \times \ln 2 \approx$
\textbf{20 millions d'années}.
\end{solution}
